
\frontslide{Organisation des fichiers}


\begin{frame}[fragile]{Enregistrements, blocs et fichiers}


Contenu de ce cours\,: la représentation \alert{physique} d'une base de données (par le SGBD).

\vfill

\begin{redbullet}
   \item les enregistrements et leur adressage
   \item le stockage des enregistrements dans des blocs
   \item la séquence des blocs formant un fichier.
\end{redbullet}


\smallskip

Une application n'accède \alert{jamais} directement à la représentation physique.

\vfill


\begin{blocgauche}
\alert{Ces diapositives correspondent au support en ligne disponible sur
le site \url{http://sys.bdpedia.fr/stfichiers.html}}
\end{blocgauche}

\end{frame}

\begin{frame}{Enregistrements}

Un enregistrement = une suite de {\em champs}
stockant les valeurs des attributs.

\vfill

\begin{small}
\begin{tabular}{l|c}
{\bf Type}  & {\bf Taille en octets}  \\
\hline
%\multicolumn{2}{|c|}{Numériques}\\
\sqlcom{INTEGER} & 4 \\
\sqlcom{FLOAT} & 4 \\
\sqlcom{DOUBLE PRECISION} & 8 \\
\sqlcom{DECIMAL} (\param{M}, \param{D}) &  \param{M}, (\param{D}+2
                                 si \param{M} $<$ \param{D}) \\
\hline
\sqlcom{CHAR}(\param{M}) & \param{M} \\
\sqlcom{VARCHAR}(\param{M}) & \param{L+1} avec \param{L} $\leq$ \param{M} \\
\hline
\end{tabular}
\end{small}


\vfill
\begin{blocgauche}
Certains champs ont une taille \alert{variable};
d'autres peuvent être à \texttt{NULL}: pas de valeur.
\end{blocgauche}
\end{frame}

\begin{frame}{Tailles variables et valeurs {\tt NULL}}

Si tous les champs avaient de taille fixe et ont
une valeur\,: pas de problème.

\vfill

En pratique\,: certains champs ont une taille variable
ou sont à {\tt NULL}.

\vfill

\begin{redbullet}
  \item pour les champs de taille variable\,: on précède
           la valeur par la taille exacte\,;

  \item  pour les valeurs {\tt NULL}\,:
    on peut indiquer une taille 0 (Oracle)\,;
      on peut créer un \guill{masque} de bits.
\end{redbullet}

\vfill
\begin{blocgauche}
La norme n'impose pas de choix, mais dans tous
les cas il faut gérer une \alert{information
complémentaire} sur les enregistrements, et
la placer dans un \alert{entête}.
\end{blocgauche}
\end{frame}

\begin{frame}{En-tête d'enregistrement}

Les informations complémentaires sont stockées
dans l'en-tête d'un enregistrement. 

\vfill
Exemple\,:  
\begin{redbullet}
  \item table \tble{Film} (id INT, titre VARCHAR(50), année INT)
   \item Enregistrement (123, 'Vertigo', NULL)
\end{redbullet}

\vfill

\figSlideWithSize{exenr}{8cm}

\vfill

\begin{blocgauche}
Le contenu de l'enregistrement se décrypte 
au moment de la lecture.
\end{blocgauche}
\end{frame}

\begin{frame}{Blocs}

\begin{redbullet}
  \item on veut qu'un enregistrement soit stocké de manière \alert{stable}
     pour avoir une adresse \alert{fixe}\,; \alert{MAIS}:
   \item on doit envisager le cas où la taille d'un enregistrement
                   varie\,;
   \item on doit envisager  une méthode pour gérer un déplacement si nécessaire.
\end{redbullet}

\vfill
Une solution (Oracle): chaque bloc gère les adresses internes de ses enregistrements.

\vfill


\figSlideWithSize{blocenr3}{8cm}

\end{frame}


\begin{frame}{Adressage des enregistrements}

Avec le schéma précédent, \alert{l'adresse} d'un enregistrement est constituée

\begin{redbullet}
  \item D'une adresse \alert{physique}, celle du bloc. Ex: F1.12
  \item D'une adresse \alert{logique}, interne au bloc. Ex: 16
\end{redbullet}
L'adresse est \alert{fixe} et utilisée pour les \alert{accès direct} (indexation).

\vfill

Avec cette solution:
\vfill

\begin{blocgauche}
\begin{redbullet}
  \item un enregistrement peut changer de taille, avec réorganisation
                interne\,;

  \item Plus de place dans le bloc? On peut déplacer l'enregistrement.
\end{redbullet}
\end{blocgauche}


\end{frame}

\begin{frame}{Réorganisation du stockage}

\alert{Un enregistrement s'agrandit, mais qu'il reste
                de la place dans le bloc}:

\begin{redbullet}
  \item On modifie l'organisation interne\,; la table \alert{locale} d'adressage est modifiée 
\end{redbullet}

\vfill
\alert{Un enregistrement s'agrandit, plus de place dans le bloc}:
\begin{redbullet}
   \item on   déplace l'enregistrement  et on crée un \alert{chaînage
                  dans l'en-tête du bloc}.
\end{redbullet}

\vfill
\begin{blocgauche}
La création de chaînage pénalise les performances
$\Rightarrow$ si possible laisser de l'espace libre 
dans un bloc.
\end{blocgauche}
\end{frame}

\begin{frame}{Exemple d'évolution avec chaînage}


\figSlideWithSize{blocenr4}{11cm}


\vfill
\begin{blocgauche}
On agrandit deux fois successivement l'enregistrement
F1.12.46.
\end{blocgauche}

\end{frame}



\begin{frame}{Fichier = séquence de blocs}

Un fichier est simplement une séquence de blocs contenant
les données d'un même objet (table, index).

\vfill

La qualité d'un fichier est conditionnée par\,:

\begin{redbullet}
   \item la bonne utilisation de l'espace (idéalement
                  tous les blocs sont pleins)\,;

  \item le stockage le plus contigu possible (même piste,
            même cylindre, etc).
   \item peu ou pas de chaînage
\end{redbullet}

\vfill

\begin{blocgauche}
Si les blocs sont très dispersés sur le disque, on
aboutit à une \alert{fragmentation} très pénalisante.
\end{blocgauche}

\end{frame}

\begin{frame}{Opérations sur un fichier}

Voir le support. Essentiellement:


\vfill

\begin{redbullet}
   \item En l'absence d'index, le seul moyen de
rechercher un enregistrement est de parcourir
\alert{séquentiellement} le fichier.

  \item Si le fichier est trié il est possible d'effectuer des recherches
   par \alert{dichotomie}.
\end{redbullet}

\vfill

\begin{blocgauche}
En général, les fichiers gérés par les SGBD sont de simple
fichiers séquentiels, sans ordre.
\end{blocgauche}

\end{frame}

\begin{frame}{Résumé}

La gestion des fichiers est entièrement à la charge du SGBD.

\vfill 
\alert{Retenir}:

\begin{redbullet}
  \item Les enregistrements sont stockés dans des blocs.
  \item Le contenu des blocs d'adapte au changement de taille d'un enregistrement après insertion.
  \item \alert{Il faut un dispositif stable d'adressage de chaque enregistrement}
  \item Les fichiers sont des séquences de blocs, les moins fragmentés possibles.
\end{redbullet}

\vfill

\begin{blocgauche}
L'adressage stable est essentiel pour permettre l'indexation
et l'accès direct aux enregistrements.
\end{blocgauche}

\end{frame}